\documentclass[10pt,a4paper]{article}
\usepackage{fontspec,amsmath,amssymb,graphicx}
\usepackage[margin=18mm]{geometry}
\setmainfont{Times New Roman}
\setlength{\parindent}{0pt}
\setlength{\parskip}{5pt}
\allowdisplaybreaks
\sloppy
\emergencystretch=2em
\begin{document}
\section*{BT\_2\_3\_15 — Lời giải bổ sung}
\subsection*{Kiến thức và trình tự giải}
\textbf{Mục tiêu.} Tận dụng phần cấp hai có dạng bình phương một đạo hàm theo hướng để giảm PDE thành ODE có tham số.

\textbf{Bước 1 — Nhận ra hướng đạo hàm.} Nếu $D=a\partial_x+b\partial_y$ có hệ số hằng, $D^2u=a^2u_{xx}+2ab u_{xy}+b^2u_{yy}$. Với hệ số biến thiên phải dùng


\[
D^2u=a^2u_{xx}+2ab u_{xy}+b^2u_{yy}+(Da)u_x+(Db)u_y.
\]


So sánh cả hạng cấp một trước khi kết luận PDE bằng $D^2u+\lambda Du+\mu u=f$.

\textbf{Bước 2 — Tìm bất biến $s$ và biến tiến $z$.} Giải $Ds=0$, chọn $Dz=k\ne0$. Kiểm tra $J=\partial(s,z)/\partial(x,y)\ne0$. Khi $k$ hằng, $D=k\partial_z$; PDE trở thành ODE theo $z$, trong đó $s$ là tham số. Nếu còn $v_s$, phải giữ lại, chỉ các trường hợp nó triệt tiêu mới giải trực tiếp như ODE.

\textbf{Bước 3 — Giải ODE.} Với $v_{zz}=g(s,z)$, tích phân hai lần cho $v=zF(s)+G(s)+\int_{z_0}^z(z-t)g(s,t)dt$. Với $a v_{zz}+b v_z+c v=g$, giải đa thức $a\lambda^2+b\lambda+c=0$, giữ hai hàm tùy ý theo $s$. Nghiệm kép có thêm hệ số $z$; có thể tìm nghiệm riêng dạng mũ hoặc đa thức.

\textbf{Bước 4 — Dùng dữ liệu đúng biến.} Tính $u_x=s_xv_s+z_xv_z$ và $u_y=s_yv_s+z_yv_z$. Nếu dữ liệu ở $s=0$, chỉ biết $F(0),G(0)$ hoặc đạo hàm của chúng; không suy ra toàn bộ các hàm. So sánh hệ số của $e^z$, $z$, $z^2$ để tìm các hằng số tương thích.

\textbf{Bước 5 — Kiểm tra tính duy nhất và nghiệm cụ thể.} Khi đường dữ liệu gặp tất cả đặc trưng một cách độc lập, giải hệ cho $F,G$. Khi đường dữ liệu là một đặc trưng, mô tả các hàm còn tùy ý và đưa hai lựa chọn rõ ràng. Cuối cùng thế lại phương trình gốc, không chỉ ODE mới.
\subsection*{Lời giải áp dụng cho bài}
\textbf{(a) Phân loại và hướng đạo hàm.} Trên $0<x<\pi/2$, $y>0$, $A=\tan^2x$, $B=-y\tan x$, $C=y^2$, $\delta=0$. Đặt $D=\tan x\partial_x-y\partial_y$. Ta có
\[D^2u=\tan^2x\,u_{xx}-2y\tan x\,u_{xy}+y^2u_{yy}
+\tan x\sec^2x\,u_x+y u_y.\]
Trừ $Du$ cho hệ số $u_x$ là $\tan x(\sec^2x-1)=\tan^3x$, hệ số $u_y$ là $2y$, đúng đề. PDE là $(D^2-D)u=\tan^3x$.

\textbf{Tọa độ đặc trưng.} Từ $dy/dx=-y/\tan x$, tích phân $dy/y=-\cot x\,dx$ cho $s=y\sin x$ không đổi. Chọn $z=\log\sin x$: $Dz=1$, $Ds=0$, $J=-\cos x\ne0$. Miền mới $s>0$, $z<0$. Dạng chính tắc $v_{zz}-v_z=\tan^3x(z)$.

\textbf{(b) Giải ODE.} Nghiệm riêng $u_p=x$ vì $Dx=\tan x$, $D^2x-Dx=\tan x\sec^2x-\tan x=\tan^3x$. Phần thuần nhất có nghiệm $1,e^z$, nên
\[\boxed{u=x+F(y\sin x)+\sin x\,G(y\sin x).}\]

\textbf{(c) Tìm $a,b$ và hai nghiệm.} Giả sử nghiệm có vết liên tục đến $y=0$ với các hàm $F,G$ đủ trơn tại $s=0$. Trên biên, $x+F(0)+\sin xG(0)=ax+b\sin x$. Ba hàm $1,x,\sin x$ độc lập trên $(0,\pi/2)$, suy ra
\[\boxed{a=1,\quad F(0)=0,\quad G(0)=b,\quad b\in\mathbb R\text{ tùy ý}.}\]
Chọn $F=0$, $G=b$ và $F(s)=s$, $G=b$ cho
\[u_1=x+b\sin x,\qquad u_2=x+b\sin x+y\sin x.\]
Kiểm tra $D(\sin x)=\sin x$, nên $(D^2-D)(b\sin x)=0$; $D(y\sin x)=0$, nên phần chênh cũng thỏa thuần nhất. Cả hai có vết $x+b\sin x$ trên $y=0$. Có vô số nghiệm; dữ liệu biên là một đặc trưng, chỉ xác định giá trị hai hàm tùy ý tại $0$.

\textbf{Chi tiết chứng minh độc lập và dữ liệu biên.} Từ $(1-a)x+F(0)+(G(0)-b)\sin x=0$, đạo hàm hai lần theo $x$ cho $-(G(0)-b)\sin x=0$. Trên khoảng đang xét $\sin x>0$, nên $G(0)=b$; phương trình còn $(1-a)x+F(0)=0$, đạo hàm cho $a=1$ rồi $F(0)=0$. Không có điều kiện nào còn lại để cố định $b$.
\begin{center}\includegraphics[width=.97\linewidth,height=11cm,keepaspectratio]{Do_thi/BT_2_3_15.png}\end{center}
\textit{Đồ thị minh họa; với nghiệm số, xem phương pháp và tham số trong mã Python kèm theo.}\end{document}